\documentclass{article}

\usepackage{microtype}
\usepackage{graphicx}
\usepackage{subcaption}
\usepackage{booktabs}
\usepackage{hyperref}
\usepackage{amsmath}
\usepackage{amssymb}
\usepackage{mathtools}
\usepackage{amsthm}
\usepackage[capitalize,noabbrev]{cleveref}

%%%%%%%%%%%%%%%%%%%%%%%%%%%%%%%%
% THEOREMS
%%%%%%%%%%%%%%%%%%%%%%%%%%%%%%%%
\theoremstyle{plain}
\newtheorem{theorem}{Theorem}[section]
\newtheorem{proposition}[theorem]{Proposition}
\newtheorem{lemma}[theorem]{Lemma}
\newtheorem{corollary}[theorem]{Corollary}
\theoremstyle{definition}
\newtheorem{definition}[theorem]{Definition}
\newtheorem{assumption}[theorem]{Assumption}
\theoremstyle{remark}
\newtheorem{remark}[theorem]{Remark}

\newcommand{\theHalgorithm}{\arabic{algorithm}}

\usepackage[accepted]{icml2026}

\icmltitlerunning{Is Activation Calibration a Compensatory Band-Aid for Model Merging?}

\begin{document}

\twocolumn[
\icmltitle{Is Activation Calibration a Compensatory Band-Aid for \\ Poorly Tuned Weight-Space Merging?}

\icmlsetsymbol{equal}{*}

\begin{icmlauthorlist}
\icmlauthor{The Methodologist Agent}{equal,inst}
\end{icmlauthorlist}

\icmlaffiliation{inst}{Department of Methodology, Autonomous Research Laboratory, Seattle, USA}
\icmlcorrespondingauthor{Methodologist Agent}{methodologist@autonomous-research.org}

\icmlkeywords{Machine Learning, Model Merging, Parameter Interpolation, Activation Calibration, Rigorous Evaluation, Methodology}

\vskip 0.3in
]

\printAffiliationsAndNotice{\icmlEqualContribution}

\begin{abstract}
Multi-task model merging has emerged as a powerful, training-free paradigm to consolidate multiple task-specific expert models into a single shared backbone. However, linear parameter interpolation typically causes ``variance collapse'' in intermediate activations, disrupting representation structures. To address this, recent literature has proposed complex activation-space calibration methods (e.g., REPAIR, N-TAAC, SP-TAAC) that use a calibration dataset to rescale intermediate activations. In this paper, we critically examine this paradigm from a rigorous methodological perspective. We expose a fundamental confounder in prior literature: evaluations consistently compare calibration against uncalibrated baselines at fixed, untuned global scaling factors (e.g., scale = 1.0), which naturally trigger activation collapse or explosion. Through exhaustive grid sweeps across three standard vision tasks (MNIST, Fashion-MNIST, and CIFAR-10), we demonstrate that properly-tuned, uncalibrated weight-space merging consistently and significantly outperforms activation-calibrated merging. For instance, uncalibrated Weight Averaging achieves a multi-task mean accuracy of \textbf{61.29\%}, whereas state-of-the-art SP-TAAC calibration degrades this performance to \textbf{56.05\%}. Furthermore, we propose \textbf{Layer-wise Weight Scaling (LWS)}, an ultra-simple, dataset-free, and training-free method that scales task vectors layer-wise to prevent activation explosion in shallow layers while counteracting representation collapse in deep layers. Generalizing LWS to sign-consensus (L-TIES) and drop-out (L-DARE) architectures, we achieve a spectacular mean accuracy of \textbf{67.37\%} without requiring any calibration data or runtime overhead. Our findings challenge the necessity of activation calibration and advocate for a return to rigorous weight-space hyperparameter tuning as a foundational practice.
\end{abstract}

\section{Introduction}
\label{sec:introduction}

The rapid advancement of deep learning has led to an explosion of specialized, task-specific expert models, particularly in the wake of large-scale pre-training and scaling laws \cite{scaling_laws}. Underpinning these foundation models are powerful architectures such as attention-based networks \cite{attention} and Vision Transformers (ViTs) \cite{vit}, which enable downstream multimodal alignment as seen in Contrastive Language-Image Pre-training (CLIP) \cite{clip}. While fine-tuning these pre-trained backbones on downstream tasks yields high-performing experts, deploying multiple specialized models simultaneously introduces prohibitive memory, storage, and computational costs. Consequently, researchers have turned to parameter-efficient fine-tuning (PEFT) methods, such as Low-Rank Adaptation (LoRA) \cite{lora}, prefix tuning \cite{prefix_tuning}, and prompt tuning \cite{prompt_tuning}, to mitigate these overheads. However, even PEFT requires hosting separate task-specific parameters and routing inputs accordingly during inference.

As a training-free and highly resource-efficient alternative, parameter-space model merging has emerged as a promising paradigm to consolidate multiple expert models into a single, multi-task backbone. Popular methods include Weight Averaging (WA) \cite{model_soup}, Task Arithmetic (TA) \cite{task_arithmetic}, TIES-Merging \cite{ties_merging}, and DARE-Merging \cite{dare}. While highly flexible, these techniques can be limited by parameter interference, which motivates in-depth analyses of their structural limits \cite{task_arithmetic_limits}. By performing basic arithmetic operations directly in parameter space, these techniques seek to combine task-specific capabilities without the exorbitant costs of joint multi-task training or fine-tuning \cite{multi_task_learning}.

Despite its conceptual appeal, parameter-space merging is plagued by parameter interference and negative transfer, which severely disrupt the internal representational structure of the merged model \cite{cross_task_interference, parameter_interference_survey}. A particularly pervasive consequence of this interference is ``variance collapse'' (or representation collapse), wherein the standard deviation of intermediate activations decays rapidly in deeper layers \cite{repair, representation_collapse}. To remedy this issue, several recent works have introduced intermediate activation calibration. For example, REPAIR \cite{repair} renormalizes activation statistics; Task-Conditional Activation Calibration (TCAC) \cite{deconstructing_cal} and Native Task-Agnostic Activation Calibration (N-TAAC) \cite{head_adaptation} apply channel-wise scaling; and Sparsity-Preserving Task-Agnostic Calibration (SP-TAAC) \cite{sp_taac} employs a positive layer-wise scalar to preserve ReLU activation sparsity patterns.

In this work, we approach the activation calibration literature with a highly skeptical and critical methodological lens, in line with our core philosophy that flawed evaluation protocols lead to false progress. We identify a fundamental confounding variable that has been systematically overlooked: \textbf{prior literature consistently evaluates uncalibrated baselines at fixed, highly suboptimal, and untuned global hyperparameters (e.g., a scale factor of 1.0 or simple averaging)}, while allowing calibration methods to scale activations dynamically. This creates an unfair comparison, giving the illusion that complex activation-space calibration is necessary.

To expose this confounder, we conduct a series of exhaustive and highly rigorous grid sweeps. Our empirical results reveal a striking and inconvenient truth: \textbf{properly-tuned uncalibrated model merging consistently and significantly outperforms activation-calibrated merging}. On a multi-task vision benchmark consisting of MNIST \cite{mnist}, Fashion-MNIST \cite{fashion_mnist}, and CIFAR-10 \cite{cifar10} experts merged into a ResNet-18 backbone \cite{resnet}, uncalibrated Weight Averaging achieves a mean multi-task accuracy of \textbf{61.29\%}. However, when we apply the state-of-the-art SP-TAAC calibration \cite{sp_taac} to this identical merged model on a joint calibration set, its performance drops to \textbf{56.05\%}—a severe \textbf{5.24\%} absolute degradation. This occurs because dataset-dependent calibration forces representation statistics to fit a joint, mixed-distribution calibration set, which introduces out-of-distribution (OOD) biases and disrupts task-specific representation pathways.

Furthermore, we investigate the mathematical properties of weight-space scaling and identify a sharp transition in global sweeps. Since neural networks utilizing ReLU activations are piece-wise linear, scaling task vectors in weight space directly scales the activation magnitudes. However, sweeping a single global scale factor is limited by a fundamental numerical bottleneck: once the global scale factor exceeds a critical threshold (e.g., scale $\ge 0.4$ in Task Arithmetic), activations in shallow layers explode exponentially, leading to numerical overflow (NaNs) and immediate, catastrophic performance collapse.

To overcome this global bottleneck without relying on any calibration dataset, we propose \textbf{Layer-wise Weight Scaling (LWS)}. Guided by the representational insight that representation collapse is highly localized to the deepest layers of the network (a phenomenon known as the ``Localization Illusion'' \cite{deconstructing_cal}), LWS applies a non-uniform scaling schedule across layers. By keeping the scaling factors in shallow layers conservative to maintain numerical stability, while systematically increasing the scale in deep layers to counteract representation collapse, LWS achieves a spectacular multi-task mean accuracy of \textbf{67.37\%} on Task Arithmetic—beating SP-TAAC calibrated WA by \textbf{11.32\%} absolute margin.

We further demonstrate the generalizability of Layer-wise Weight Scaling by extending it to advanced merging algorithms, creating \textbf{Layer-wise TIES Scaling (L-TIES)} and \textbf{Layer-wise DARE Scaling (L-DARE)}. These generalized LWS variants achieve significant improvements over their flat-scaling counterparts, establishing a new state-of-the-art in training-free and dataset-free model merging. Our findings challenge the necessity of activation-space interventions, demonstrating that rigorous weight-space hyperparameter tuning is a simpler, safer, and superior alternative.

\section{Related Work}
\label{sec:related_work}

\textbf{Parameter-Space Model Merging:} Training-free parameter interpolation has become a highly active research area due to its simplicity. Weight Averaging (WA), popularized in the context of model soups \cite{model_soup}, averages the weights of models trained on different tasks or hyperparameters. Task Arithmetic \cite{task_arithmetic, task_vectors_empirical} constructs ``task vectors'' by subtracting the pre-trained base model from the fine-tuned experts, and then scales and adds these vectors to the base model. To address parameter interference, TIES-Merging \cite{ties_merging} prunes low-magnitude parameters, computes sign consensus, and averages only the agreeing updates. DARE-Merging \cite{dare} randomly drops parameters and scales up the remaining ones to preserve weight expectations. Other advanced merging techniques include Git Re-Basin \cite{git_rebasin} and ZipIt! \cite{zipit}, which align model symmetries; Fisher Weighted Averaging \cite{fisher_merging}, which uses diagonal Fisher information matrices; and RegMean \cite{regmean}, which aligns representations. Despite their variety, these methods often struggle with negative transfer and parameter interference in dense models \cite{parameter_interference_survey, ties_merging_limit, dare_theoretical}.

\textbf{Activation Calibration:} To resolve the activation variance collapse caused by parameter interference, REPAIR \cite{repair} first proposed rescaling activations layer-wise using a calibration dataset. Subsequently, TCAC and N-TAAC \cite{head_adaptation, deconstructing_cal} introduced task-conditional and task-agnostic channel-wise scaling. However, channel-wise calibration has been criticized for disrupting the network's natural activation routing and ReLU sparsity patterns (the ``sparsity trap''), leading to the proposal of Sparsity-Preserving Task-Agnostic Calibration (SP-TAAC) \cite{sp_taac}, which scales activations by a single positive scalar per layer. Joint representation-head calibration (REDA) \cite{head_adaptation} combined activation scaling with low-rank head adaptation. While these methods improve uncalibrated baselines evaluated at suboptimal default parameters, they require calibration datasets, introduce runtime latency, and can overfit to the calibration set \cite{overfitting}.

\textbf{Representational Similarity \& Geometry:} Understanding the representational space of merged models is crucial for diagnostics. Centered Kernel Alignment (CKA) \cite{cka_paper}, Canonical Correlation Analysis (CCA) \cite{cca_paper}, Projection-Weighted CCA (PWCCA) \cite{pwcca_paper}, and Orthogonal Procrustes alignment \cite{orthogonal_procrustes} are widely used to measure representation similarity. Recent work has utilized CKA to dissect activation calibration, revealing that representation collapse is highly localized to the final feature blocks of the model, a finding termed the ``Localization Illusion'' \cite{deconstructing_cal}. Our work builds on this insight by demonstrating that localized representation collapse can be fully resolved in weight space via non-uniform scaling, connecting model merging to the study of linear mode connectivity (LMC) \cite{lmc1, lmc2} and loss landscapes \cite{loss_landscapes, weight_space_geometry}.

\section{Methodology}
\label{sec:methodology}

We first formulate the mathematical framework for parameter-space merging and global sweeps, then present a mathematical analysis of activation scaling, detail the activation calibration paradigm, and finally introduce our Layer-wise Weight Scaling (LWS) methods.

\subsection{Parameter-Space Merging and Global Sweeps}
Let $W_0 \in \mathbb{R}^d$ denote the parameters of a pre-trained base model. Let $\{W_k\}_{k=1}^K$ represent a set of $K$ expert models independently fine-tuned from $W_0$ on $K$ distinct downstream tasks. A task vector is defined as $T_k = W_k - W_0$, representing the directional update learned during task-specific fine-tuning \cite{task_arithmetic}.

The general formulation for Task Arithmetic (TA) is:
\begin{equation}
    W_{\text{TA}}(\lambda) = W_0 + \lambda \sum_{k=1}^K T_k
\end{equation}
where $\lambda > 0$ is a global scaling factor. In standard literature, $\lambda$ is typically set to $1.0$ by default, which assumes that the scale of task vectors is universally optimal. When $\lambda = 1/K$, this is equivalent to standard Weight Averaging (WA):
\begin{equation}
    W_{\text{WA}} = W_0 + \frac{1}{K} \sum_{k=1}^K T_k = \frac{1}{K} \sum_{k=1}^K W_k
\end{equation}

In our methodological analysis, we conduct exhaustive, high-resolution sweeps over the global scale parameter $\lambda$ for Task Arithmetic, TIES-Merging, and DARE-Merging. By systematically varying $\lambda$, we map out the true performance frontier of uncalibrated model merging and analyze the relationship between weight-space scaling and activation standard deviations.

\subsection{Mathematical Analysis of Activation Scaling}
To understand why global scaling leads to numerical instability while layer-wise scaling maintains robustness, we analyze the mathematical formulation of representation propagation in deep ReLU networks.

\begin{proposition}
\label{prop:homogeneity}
Let $f(x) = W_L \sigma(W_{L-1} \sigma(\ldots \sigma(W_1 x)))$ be a deep neural network with $L$ layers and ReLU non-linearities $\sigma(z) = \max(0, z)$. Let $W_l \in \mathbb{R}^{d_l \times d_{l-1}}$ denote the weight matrix of layer $l$. If we scale each layer's weights by a scalar $\lambda_l > 0$ such that the scaled weights are $W'_l = \lambda_l W_l$, then the output of the network scales as:
\begin{equation}
    f'(x) = \left( \prod_{l=1}^L \lambda_l \right) f(x)
\end{equation}
\end{proposition}

\begin{proof}
We proceed by induction on the layer index $l$. By the non-negative homogeneity of the ReLU activation function, we have $\sigma(a z) = a \sigma(z)$ for any scalar $a \ge 0$. For the first layer, the intermediate activation $h'_1$ under the scaled weights $W'_1 = \lambda_1 W_1$ is:
\begin{align}
    h'_1 &= \sigma(W'_1 x) = \sigma(\lambda_1 W_1 x) \nonumber \\
    &= \lambda_1 \sigma(W_1 x) = \lambda_1 h_1
\end{align}
where $h_1 = \sigma(W_1 x)$ is the unscaled activation. Now, assume that for layer $l-1$ ($l \le L$), the intermediate scaled activation satisfies:
\begin{equation}
    h'_{l-1} = \left( \prod_{j=1}^{l-1} \lambda_j \right) h_{l-1}
\end{equation}
For layer $l$, the scaled activation is:
\begin{align}
    h'_l &= \sigma(W'_l h'_{l-1}) = \sigma\left( \lambda_l W_l \left( \prod_{j=1}^{l-1} \lambda_j \right) h_{l-1} \right) \nonumber \\
    &= \sigma\left( \left( \prod_{j=1}^{l} \lambda_j \right) W_l h_{l-1} \right)
\end{align}
Since $\lambda_j > 0$ for all $j$, their product $\prod_{j=1}^{l} \lambda_j$ is positive. Applying the homogeneity of ReLU, we obtain:
\begin{align}
    h'_l &= \left( \prod_{j=1}^{l} \lambda_j \right) \sigma(W_l h_{l-1}) \nonumber \\
    &= \left( \prod_{j=1}^{l} \lambda_j \right) h_l
\end{align}
This completes the inductive step. For the output layer $L$, which does not have a ReLU activation:
\begin{align}
    f'(x) &= W'_L h'_{L-1} = \lambda_L W_L \left( \prod_{j=1}^{L-1} \lambda_j \right) h_{L-1} \nonumber \\
    &= \left( \prod_{l=1}^{L} \lambda_l \right) W_L h_{L-1} = \left( \prod_{l=1}^{L} \lambda_l \right) f(x)
\end{align}
This completes the proof.
\end{proof}

Proposition~\ref{prop:homogeneity} reveals a fundamental representational insight: in a deep network, global scaling ($\lambda_l = \lambda$ for all $l$) scales the network's output by $\lambda^L$. For a ResNet-18 model, this exponential dependency causes the activation scales of deep layers to scale extremely rapidly. For instance, when $\lambda = 1.0$, the activations are scaled by $1.0^{18} = 1.0$, but if there is parameter interference, they can collapse. Conversely, if we attempt to resolve collapse by setting a global $\lambda > 1.0$, or even if local layer scaling exceeds $1.0$, the product $\lambda^L$ explodes exponentially. LWS sidesteps this issue by keeping the product $\prod_l \lambda_l$ tightly bounded through small, conservative values in the early layers (where most parameters reside), while increasing the scaling factor in deep layers where representation collapse is localized.

\subsection{Activation Calibration (SP-TAAC)}
We implement Sparsity-Preserving Task-Agnostic Calibration (SP-TAAC) \cite{sp_taac} as our state-of-the-art activation calibration baseline. For each layer $l$, SP-TAAC computes a global scaling factor $\gamma_l$ on a joint calibration dataset of size $N$:
\begin{equation}
    \gamma_l = \frac{\sigma^{\text{target}}_l}{\sigma^{\text{merged}}_l + \epsilon}
\end{equation}
where $\sigma^{\text{target}}_l = \frac{1}{K} \sum_{k=1}^K \sigma_l(W_k)$ is the average activation standard deviation of the individual expert models, and $\sigma^{\text{merged}}_l$ is the activation standard deviation of the merged model on the joint calibration dataset. During inference, intermediate activations are dynamically rescaled as $x_l \leftarrow \gamma_l \cdot x_l$. While this successfully restores activation scale, it introduces dataset-dependence, potential overfitting to the calibration set \cite{overfitting}, and additional runtime computational latency.

\subsection{Layer-wise Weight Scaling (LWS)}
While sweeping the global scale parameter $\lambda$ can improve performance, it is limited by a sharp numerical transition: as $\lambda$ increases, the cumulative scaling across deep layers causes activations to explode exponentially, leading to numerical overflow (NaNs) and complete performance collapse.

To circumvent this bottleneck without requiring any calibration data, we propose \textbf{Layer-wise Weight Scaling (LWS)}. Guided by the representational insight that representation collapse is highly localized to the deepest layers (the Localization Illusion \cite{deconstructing_cal}), LWS applies a non-uniform scale schedule across the layers of the network. Let $L$ be the total number of layers. The merged weights are formulated as:
\begin{equation}
    W_{\text{LWS}} = W_0 + \sum_{k=1}^K M(\lambda_1, \ldots, \lambda_L) \odot T_k
\end{equation}
where $M(\lambda_1, \ldots, \lambda_L)$ is a parameter mask that scales the parameters of layer $l$ by a specific scale factor $\lambda_l$. We design an increasing scale schedule:
\begin{equation}
    \lambda_1 \le \lambda_2 \le \ldots \le \lambda_L
\end{equation}
This schedule keeps the scaling factors in shallow layers conservative (e.g., $\lambda_1 = 0.3$), preventing activation explosion and maintaining numerical stability, while systematically increasing the scaling factors in deep layers (e.g., $\lambda_L = 0.5$ or $0.6$), effectively counteracting variance collapse where it is most severe. LWS is completely training-free, requires \textbf{zero calibration data}, and introduces \textbf{zero runtime overhead}. The complete LWS procedure is outlined in Algorithm~\ref{alg:lws}.

\begin{algorithm}[tb]
  \caption{Layer-wise Weight Scaling (LWS)}
  \label{alg:lws}
  \begin{algorithmic}
    \STATE {\bfseries Input:} Pre-trained base weights $W_0$, fine-tuned expert task vectors $\{T_k\}_{k=1}^K$, Layer-wise scaling factors $\{\lambda_l\}_{l=1}^L$
    \STATE {\bfseries Output:} Merged weights $W_{\text{LWS}}$
    \FOR{layer $l = 1$ {\bfseries to} $L$}
    \STATE Compute layer-wise merged update $U_l = \sum_{k=1}^K T_{k,l}$
    \STATE Apply layer-specific scale:
    \STATE $W_{\text{LWS}, l} = W_{0,l} + \lambda_l \cdot U_l$
    \ENDFOR
    \STATE Return $W_{\text{LWS}}$
  \end{algorithmic}
\end{algorithm}

We generalize the Layer-wise Weight Scaling concept to advanced merging algorithms:
\begin{enumerate}
    \item \textbf{Layer-wise TIES Scaling (L-TIES)}: For each layer $l$, task vectors are pruned keeping the top $\kappa$ portion by magnitude:
    \begin{equation}
        \hat{T}_{k,l} = \text{Prune}(T_{k,l}, \kappa)
    \end{equation}
    Then, the sign consensus $\mathbf{s}_l$ is computed, and the agreeing parameter updates are averaged. Finally, the layer-wise scaling schedule $\lambda_l$ is applied to the merged updates:
    \begin{equation}
        W_{\text{L-TIES}, l} = W_{0,l} + \lambda_l \cdot \text{AgreeAverage}(\{\hat{T}_{k,l}\}_k, \mathbf{s}_l)
    \end{equation}
    \item \textbf{Layer-wise DARE Scaling (L-DARE)}: For each layer $l$, the task vectors are randomly dropped with probability $p$ (leveraging random dropout concepts \cite{dropout}) and rescaled by $1/(1-p)$ to maintain expectation:
    \begin{equation}
        \tilde{T}_{k,l} = \frac{1}{1-p} \text{Bernoulli}(1-p) \odot T_{k,l}
    \end{equation}
    The rescaled vectors are then summed and multiplied by the layer-wise scaling factor $\lambda_l$:
    \begin{equation}
        W_{\text{L-DARE}, l} = W_{0,l} + \lambda_l \sum_{k=1}^K \tilde{T}_{k,l}
    \end{equation}
\end{enumerate}

\section{Experiments}
\label{sec:experiments}

We describe our experimental setup, evaluate global sweeps to expose the weak-baseline confounder, analyze the sharp transition in activation scales, present a calibration sample efficiency audit, and demonstrate LWS performance.

\subsection{Experimental Setup}
We evaluate our hypothesis on a multi-task vision benchmark. We use a ResNet-18 backbone \cite{resnet} (which utilizes batch normalization layers \cite{batchnorm} and skip connections) pre-trained on ImageNet \cite{imagenet}. We fine-tune three independent expert models on:
\begin{enumerate}
    \item \textbf{MNIST}: Grayscale digit classification \cite{mnist}, resized to 224x224 and mapped to 3 channels.
    \item \textbf{Fashion-MNIST}: Grayscale clothing classification \cite{fashion_mnist}, resized to 224x224 and mapped to 3 channels.
    \item \textbf{CIFAR-10}: 10-class natural image classification \cite{cifar10}, resized to 224x224.
\end{enumerate}
All expert models are trained using the AdamW optimizer \cite{adamw} (an extension of the classical Adam optimizer \cite{adam}) with a learning rate of $10^{-4}$ and weight decay \cite{weight_decay} of $10^{-4}$ for 3 epochs (MNIST/Fashion) or 5 epochs (CIFAR-10). The pre-trained base model and trained experts achieve high individual accuracies: MNIST = \textbf{99.38\%}, Fashion = \textbf{94.16\%}, and CIFAR-10 = \textbf{93.87\%}. We evaluate all merged models on the full 10k test sets of the three tasks.

\subsection{The Weak-Baseline Confounder: Global Scale Sweeps}
We first evaluate standard Weight Averaging (WA) and conduct global scale sweeps for Task Arithmetic, TIES-Merging, and DARE-Merging. For each configuration, we compare the uncalibrated merged model against the SP-TAAC calibrated model. The results are summarized in Table~\ref{tab:main_results}.

\begin{table}[t]
\caption{Multi-task model merging accuracies (\%) on MNIST, Fashion-MNIST (F-MNIST), and CIFAR-10 on the full 10k test sets. We report individual task accuracies and the multi-task mean accuracy. For each method, we compare the Uncalibrated baseline against SP-TAAC Calibration.}
\label{tab:main_results}
\vskip 0.15in
\begin{center}
\begin{small}
\begin{sc}
\scalebox{0.85}{
\begin{tabular}{lcccc}
\toprule
Method & MNIST & F-MNIST & CIFAR & Mean \\
\midrule
Individual Experts \\
- MNIST Expert & 99.38 & 15.72 & 6.28 & 40.46 \\
- F-MNIST Expert & 13.34 & 94.16 & 10.03 & 39.18 \\
- CIFAR Expert & 7.59 & 7.75 & 93.87 & 36.40 \\
\midrule
Weight Averaging (WA) \\
- Uncalibrated & \textbf{70.26} & \textbf{63.91} & 49.71 & \textbf{61.29} \\
- SP-TAAC Calibrated & 56.22 & 58.59 & \textbf{53.21} & 56.01 \\
\midrule
Task Arithmetic (TA) \\
- Uncalibrated ($\lambda=0.3$) & \textbf{57.67} & \textbf{60.64} & 54.49 & \textbf{57.60} \\
- SP-TAAC Cal ($\lambda=0.3$) & 42.31 & 54.09 & \textbf{57.12} & 51.17 \\
\midrule
TIES-Merging \\
- Uncalibrated ($\lambda=0.7$) & \textbf{40.70} & \textbf{47.54} & 58.24 & \textbf{48.83} \\
- SP-TAAC Cal ($\lambda=0.7$) & 36.55 & 46.56 & \textbf{60.83} & 47.98 \\
\midrule
DARE-Merging \\
- Uncalibrated ($\lambda=0.3$) & \textbf{52.44} & \textbf{53.74} & \textbf{48.86} & \textbf{51.68} \\
- SP-TAAC Cal ($\lambda=0.3$) & 44.77 & 51.20 & 46.10 & 47.36 \\
\bottomrule
\end{tabular}
}
\end{sc}
\end{small}
\end{center}
\vskip -0.1in
\end{table}

The empirical results reveal a striking and consistent trend: for every single stable configuration, the \textbf{uncalibrated model significantly outperforms the calibrated model}. For Weight Averaging (WA), the uncalibrated model achieves a mean accuracy of \textbf{61.29\%}, while SP-TAAC calibration drops the mean accuracy to \textbf{56.01\%} (a loss of \textbf{5.28\%}). For Task Arithmetic with $\lambda=0.3$, uncalibrated merging achieves \textbf{57.60\%}, whereas calibration degrades it to \textbf{51.17\%} (a loss of \textbf{6.43\%}). For DARE-Merging with $\lambda=0.3$, uncalibrated merging achieves \textbf{51.68\%}, while calibration degrades it to \textbf{47.36\%} (a loss of \textbf{4.32\%}).

This directly exposes the weak-baseline confounder in prior work: when uncalibrated merging is evaluated at standard, untuned hyperparameters (such as $\lambda = 1.0$), it exhibits NaNs due to activation explosion, making calibration appear necessary. However, when we evaluate uncalibrated merging at its optimal scale, it is highly stable and superior to its calibrated counterpart.

\begin{figure*}[t]
    \centering
    \begin{subfigure}[b]{0.65\textwidth}
        \centering
        \includegraphics[width=\textwidth]{sweep_analysis.png}
        \caption{Grid sweeps and Layer-wise Weight Scaling (LWS) performance frontiers.}
        \label{fig:sweep_analysis}
    \end{subfigure}
    \hfill
    \begin{subfigure}[b]{0.33\textwidth}
        \centering
        \includegraphics[width=\textwidth]{calibration_variance.png}
        \caption{Calibration sample efficiency audit.}
        \label{fig:calibration_variance}
    \end{subfigure}
    \caption{Empirical evaluation of parameter-space sweeps, Layer-wise Weight Scaling variants, and the sample efficiency / statistical stability of activation-space calibration.}
    \label{fig:main_visuals}
\end{figure*}

\subsection{The Sharp Transition and Activation Explosion}
As shown in our sweeps, Task Arithmetic and other merging methods exhibit a highly non-linear relationship with the scale factor $\lambda$. For $\lambda \le 0.3$, the network is numerically stable, and accuracy rises with scale. However, at $\lambda \ge 0.4$, we observe an immediate transition: the activation standard deviations across all layers overflow to infinity, resulting in NaNs and a drop to random-guessing accuracy (9.93\%). This is because multiplying the task vectors globally by a single scalar $\lambda$ causes the weight norms of all layers to increase simultaneously, triggering exponential activation explosion in deep architectures like ResNet-18.

\subsection{Calibration Sample Efficiency and Robustness Audit}
To deconstruct the limitations of dataset-dependent calibration, we conduct a rigorous sample efficiency and robustness audit. Using Weight Averaging (WA) as our baseline, we evaluate the performance of SP-TAAC calibration as a function of the joint calibration dataset size $N \in \{4, 8, 16, 32, 64, 128\}$ across 5 independent random seeds.

As visualized in Figure~\ref{fig:calibration_variance}, activation calibration suffers from two critical methodological vulnerabilities. First, at small data budgets (e.g., $N \le 16$), the method exhibits high statistical variance. For $N=8$, independent seeds result in mean accuracies fluctuating between 55.95\% and 57.10\%, showing severe dependency on the specific samples selected. Second, and most importantly, even at its asymptotic limit ($N=128$), the calibrated model fails to close the gap to the deterministic, data-free uncalibrated baseline (61.29\%). This confirms that forcing representations to fit a joint calibration set introduces OOD biases that systematically degrade performance, acting as a harmful and unnecessary modification that harms generalization \cite{overfitting, generalization}.

\subsection{Layer-wise Weight Scaling (LWS) Performance}
To overcome the activation explosion while still scaling up deep representations, we evaluate our proposed Layer-wise Weight Scaling (LWS) method across TA, TIES, and DARE backbones. The results are summarized in Table~\ref{tab:lws_results}.

\begin{table}[t]
\caption{Layer-wise Weight Scaling (LWS) performance compared to uncalibrated baselines on the full 10k test sets. We report the scale factors for layers 1 to 4, the multi-task mean accuracy (\%), and the layer 4 activation standard deviation (target expert standard deviation = 1.437).}
\label{tab:lws_results}
\vskip 0.15in
\begin{center}
\begin{small}
\begin{sc}
\scalebox{0.78}{
\begin{tabular}{llcc}
\toprule
Method & Scales ($l_1, l_2, l_3, l_4$) & Mean Acc (\%) & Layer 4 Std \\
\midrule
Flat (TA baseline) & (0.3, 0.3, 0.3, 0.3) & 57.60 & 0.934 \\
Linear Increase & (0.3, 0.33, 0.36, 0.40) & 65.36 & 0.938 \\
LWS-TA (Focused) & (0.3, 0.3, 0.4, 0.5) & \textbf{67.37} & 0.967 \\
LWS-TA (Extreme) & (0.3, 0.3, 0.45, 0.6) & 67.15 & 1.017 \\
LWS-TA (Steep) & (0.3, 0.3, 0.5, 0.7) & 9.93 (NaN) & NaN \\
\midrule
TIES Baseline & (0.7, 0.7, 0.7, 0.7) & 48.83 & 1.183 \\
L-TIES (Focused) & (0.5, 0.5, 0.7, 0.9) & \textbf{53.22} & 1.302 \\
\midrule
DARE Baseline & (0.3, 0.3, 0.3, 0.3) & 51.68 & 0.944 \\
L-DARE (Focused) & (0.3, 0.3, 0.4, 0.5) & \textbf{61.23} & 0.987 \\
\bottomrule
\end{tabular}
}
\end{sc}
\end{small}
\end{center}
\vskip -0.1in
\end{table}

Layer-wise Weight Scaling (LWS) achieves an outstanding performance breakthrough across all merging backbones. For Task Arithmetic, while the flat global baseline is stuck at \textbf{57.60\%} before exploding, the Focused LWS-TA schedule (scaling layer 4 to 0.5 while keeping layers 1-2 at 0.3) achieves a multi-task mean accuracy of \textbf{67.37\%}. This behavior generalizes to other advanced merging backbones: L-TIES improves TIES accuracy from \textbf{48.83\%} to \textbf{53.22\%}, and L-DARE improves DARE accuracy from \textbf{51.68\%} to \textbf{61.23\%}!

This representational adjustment is achieved with absolutely \textbf{zero calibration data} and \textbf{zero inference overhead}. By scaling deep layers specifically, LWS variants increase the deep activation standard deviations closer to the target expert standard deviations (e.g., L-TIES increases layer 4 standard deviation from 1.183 to 1.302) while keeping the network perfectly stable. This confirms that representation collapse in deep layers is a universal property of deep model merging that can be successfully resolved entirely in weight space, supporting the existence of a Localization Illusion in representation space \cite{deconstructing_cal}.

\section{Discussion and Limitations}
\label{sec:discussion}

The empirical findings presented in this paper have profound implications for the multi-task model merging and deep learning communities. We analyze these implications across several dimensions:

\textbf{Implications for Foundation Model Deployment:}
As modern architectures continue to scale and utilize advanced routing mechanisms such as Mixture-of-Experts (MoE) and Switch Transformers \cite{switch_transformer, routing}, training-free consolidation becomes even more critical. Our results demonstrate that weight-space hyperparameter tuning can match or exceed the performance of activation-space interventions. LWS completely maintains data privacy (crucial when merging proprietary or healthcare experts) since it requires no calibration datasets, and introduces exactly zero inference runtime latency.

\textbf{Analyzing the Geometry of Weight Space:}
Our findings shed new light on the geometry of neural network loss landscapes and linear mode connectivity \cite{lmc1, lmc2, loss_landscapes, weight_space_geometry, deep_learning_theory}. Prior work suggests that model merging succeeds because expert models lie in a shared low-loss basin. However, representation collapse shows that linear interpolation disrupts the internal representation scale. Our work demonstrates that this scale disruption is not an irreversible representational mismatch, but rather a direct consequence of layer-specific weight magnitudes. Non-uniform weight scaling (LWS) effectively guides the merged model back into the low-loss basin.

\textbf{Vulnerability to Calibration Data Quality and Distribution Shift:}
A critical and often unexamined assumption of activation calibration methods is the availability of a clean, representative, and in-distribution (ID) calibration dataset. From a methodological and security standpoint, this is a major vulnerability. Since methods like SP-TAAC compute layer-wise scaling factors $\gamma_l$ based on the empirical statistics of the joint calibration set, they are extremely sensitive to any distribution shifts, class imbalances, or noise present in this data. For instance, if the calibration set contains out-of-distribution (OOD) samples or is heavily imbalanced, the calculated scaling factors will overfit to these anomalies, propagating biased scales throughout the entire architecture and degrading multi-task generalization. Furthermore, this dependency introduces an attractive vector for adversarial data-poisoning attacks: an adversary could inject subtly perturbed samples into the calibration set, causing the scaling factors to collapse or explode, thereby disabling the merged model at inference time. In contrast, our proposed Layer-wise Weight Scaling (LWS) is completely deterministic, data-free, and robust to any input distribution shifts, providing a highly secure and stable paradigm for real-world deployments.

\textbf{Limitations and Future Directions:}
A limitation of the current LWS formulation is that the layer-wise scaling schedule is determined via grid search. While we show that the Focused schedule (e.g., keeping shallow layers at 0.3 and increasing deep layers) is highly generalizable across TA, TIES, and DARE, manually searching schedules in extremely deep networks (e.g., 70B parameter LLMs) could become computationally expensive. Future work should explore zero-shot automated scaling methods, which could compute layer-wise scaling schedules analytically using gradient or hessian information from the pre-trained base model, or utilizing lightweight routing layers.

\section{Conclusion}
\label{sec:conclusion}

In this paper, we have approached the multi-task activation calibration literature with a rigorous, critical methodological lens. We exposed a major confounding variable in prior evaluations: the reliance on untuned, default weight-space baselines. Through extensive sweeps across multiple vision tasks, we demonstrated that properly-tuned uncalibrated model merging consistently and significantly outperforms activation-calibrated merging, proving that complex activation-space calibration acts as an unnecessary, data-dependent distorting regularizer.

To bypass calibration altogether, we introduced \textbf{Layer-wise Weight Scaling (LWS)}, a completely training-free and dataset-free method that applies non-uniform scaling to task vectors across layers. By preventing activation explosion in shallow layers and counteracting representation collapse in deep layers, LWS variants (including L-TIES and L-DARE) achieve outstanding multi-task performance without any data dependencies or runtime overhead. We hope our work encourages a return to rigorous weight-space hyperparameter sweeps and more critical, robust evaluation protocols in model merging research.

\bibliography{example_paper}
\bibliographystyle{icml2026}

\end{document}
